\section{Methodology}
\label{sec:methodology}

\subsection{Obfuscation Testing Protocol}

The core innovation of our methodology is temporal obfuscation---systematically removing temporal and contextual information while preserving structural market mechanics. The protocol involves five transformations:

\begin{enumerate}
\item Convert absolute dates (``2024-01-16'') to relative format (``Day T+0'')
\item Replace ticker symbols (``SPY'') with generic identifiers (``INDEX\_1'')
\item Remove market events (``Fed meeting'', ``earnings'')
\item Strip volatility regime context (``VIX at 14'')
\item Eliminate day-of-week patterns (``Monday'', ``OpEx Friday'')
\end{enumerate}

Preserved information includes: net gamma exposure values, directional gamma (calls vs.\ puts), spot price and flip point levels, relative time progression, and concentration metrics. Figure~\ref{fig:obfuscation} illustrates this transformation.

\begin{figure}[H]
\centering
\includegraphics[width=\textwidth]{figures/fig01_obfuscation.png}
\caption{Temporal obfuscation process. Calendar dates, ticker symbols, and contextual markers are removed while preserving structural quantitative relationships. The transformation converts identifiable market data into anonymized numerical sequences.}
\label{fig:obfuscation}
\end{figure}

As a concrete example, the transformation converts raw market data:

\begin{small}
\begin{verbatim}
SPY, 2024-03-15: Net GEX: -$32.9B, Flip: $485.00
SPY, 2024-03-18: Net GEX: -$28.4B, Flip: $487.50
\end{verbatim}
\end{small}

\noindent into obfuscated format:

\begin{small}
\begin{verbatim}
Day T+0, INDEX_1: Net GEX: -$32.9B, Flip: $485.00
Day T+1, INDEX_1: Net GEX: -$28.4B, Flip: $487.50
\end{verbatim}
\end{small}

The quantitative values are identical; only contextual identifiers change. A model memorizing ``January 2024 SPY negative gamma'' cannot succeed when presented with ``Day T+0 INDEX\_1 GEX: $-$\$32B''. Only genuine structural understanding enables detection under these conditions.

\subsection{Causal Framework: WHO$\rightarrow$WHOM$\rightarrow$WHAT}
\label{subsec:causal}

Beyond detecting patterns, we require models to articulate causal mechanisms. This WHO$\rightarrow$WHOM$\rightarrow$WHAT framework, introduced in our single-day validation study \citep{regan2025obfuscation} and extended here to multi-day regimes, structures the model's reasoning around three components:

\textbf{WHO}: The economic actor with structural constraints (e.g., ``dealers with negative gamma exposure'').

\textbf{WHOM}: The affected market participants (e.g., ``directional traders and market makers'').

\textbf{WHAT}: The forced action and mechanism (e.g., ``must sell rallies and buy dips to maintain delta neutrality, amplifying volatility'').

Grounded in agent-based market microstructure theory \citep{grossman1988liquidity,garleanu2009demand}, this framework prevents vague pattern claims by requiring specific mechanical understanding. It implements a structured form of chain-of-thought prompting \citep{wei2022chain}, ensuring the model derives outcomes from mechanical constraints rather than recalling historical correlations.

A valid causal chain demonstrates the complete mechanism. For example: WHO---``Dealers accumulate net short gamma exposure of \$28B concentrated at the 520 strike.'' WHOM---``Market makers holding these positions are forced to continuously delta-hedge to maintain regulatory neutrality.'' WHAT---``As price rises toward 520, dealers must buy shares to offset negative delta, amplifying the upward move; as price falls, they must sell, amplifying downward moves. This creates a pro-cyclical feedback loop that persists until gamma exposure dissipates.''

\subsection{Single-Day Pattern Detection}
\label{subsec:single_day_method}

To ensure robustness, we test three narrative framings of the same underlying dealer constraint (Table~\ref{tab:patterns}). All reflect the same structural reality---dealers hedging option exposure---but use different conceptual lenses. Consistent detection across framings validates genuine understanding rather than keyword matching.

\begin{table}[H]
\centering
\caption{Three pattern framings of dealer hedging constraints. Each framing describes the same mechanical reality through a different conceptual lens.}
\label{tab:patterns}
\small
\begin{tabular}{@{}lccc@{}}
\toprule
\textbf{Pattern} & \textbf{Framing} & \textbf{WHO} & \textbf{WHAT} \\
\midrule
Gamma Positioning & Technical/Greek & Dealers with $-\Gamma$ & Pro-cyclical hedging \\
Stock Pinning & Behavioral & Market makers & Strike convergence \\
0DTE Hedging & Temporal & Option writers & Rapid rebalancing \\
\bottomrule
\end{tabular}
\end{table}

We establish three validation levels: (1) detection rate $\geq$60\% (exceeds random chance), (2) prediction accuracy $\geq$80\% (forward returns match), and (3) correct WHO$\rightarrow$WHOM$\rightarrow$WHAT specification $\geq$90\%.

\subsection{30-Day Regime Detection Framework}
\label{subsec:regime_method}

Extending from single-day patterns, we define \textit{persistent regimes} through three structural criteria applied to rolling 30-day windows. Each window contains 30 consecutive trading days of GEX values, generated with stride 1 (producing $N - 29$ overlapping windows from $N$ trading days). Windows are labeled ``Day T$-$29'' through ``Day T+0'' to preserve relative temporal ordering while removing absolute date information. For each year, this generates approximately 220 windows from $\sim$250 trading days, yielding the 1,412 real windows across 2020--2025 used in our analysis.

We classify each window using three structural criteria:

\textbf{Persistence}: Fraction of days sharing the dominant GEX sign.
\begin{equation}
P = \frac{\text{days with dominant sign}}{30} \geq 0.70
\end{equation}

\textbf{Magnitude}: Average absolute gamma exposure.
\begin{equation}
M = \frac{1}{30}\sum_{i=1}^{30} |\text{GEX}_i| \geq \$5\text{B}
\end{equation}

\textbf{Stability}: Maximum sign flips across 30 days.
\begin{equation}
S = \text{count}(\text{sign}(\text{GEX}_{i+1}) \neq \text{sign}(\text{GEX}_i)) \leq 5
\end{equation}

Windows meeting all three criteria are classified as persistent regimes (positive or negative); others are transitional or low-conviction. The 70\% persistence threshold ensures detected regimes exceed random binomial variation ($\approx 2.2\sigma$), while the \$5B magnitude threshold reflects economically significant dealer positioning. Figure~\ref{fig:regime_window} illustrates an example persistent regime.

\begin{figure}[H]
\centering
\includegraphics[width=\textwidth]{figures/fig02_regime_window.png}
\caption{30-day persistent negative regime example. Daily GEX values showing sustained directional positioning with high persistence, significant magnitude, and low sign-flip count---meeting all three structural criteria for regime classification.}
\label{fig:regime_window}
\end{figure}

\subsection{GEX Calculation}

Following standard market microstructure practice \citep{anderegg2022impact,spotgamma2021}, we calculate dealer gamma exposure from end-of-day open interest:

\begin{equation}
\text{GEX} = \sum_i \pm \text{OI}_i \times \Gamma_i \times S^2 \times 0.01 \times 100
\end{equation}

\noindent where $\Gamma_i$ is the Black-Scholes gamma \citep{black1973pricing} at strike $i$, $\text{OI}_i$ is open interest, $S$ is the underlying spot price, and the $S^2$ factor scales gamma to dollar exposure (since gamma measures $\partial^2 V / \partial S^2$, multiplying by $S^2$ produces dollar-denominated sensitivity). The sign convention reflects the dealer counterparty relationship: calls contribute positive GEX (dealers are typically long gamma on calls) and puts contribute negative GEX (dealers are typically short gamma on puts) \citep{anderegg2022impact}.

This summation yields a single scalar---\textit{net GEX}---that captures the aggregate directional bias of dealer positioning. However, the per-strike components $\text{OI}_i \times \Gamma_i$ preserve richer structural information: the spatial distribution of open interest across strikes, implied volatility skew, and localized gamma concentrations at specific price levels. Our raw chain validation (Section~\ref{subsec:raw_chain}) provides the LLM with these per-strike inputs rather than the scalar aggregate, testing whether the model can reconstruct dealer positioning from the disaggregated surface. The 30.8 percentage point improvement over the net GEX baseline (Section~\ref{sec:single_day}) suggests that scalar aggregation discards structurally informative strike-level signal.

The resulting aggregate GEX value has a direct mechanical interpretation. \textbf{Positive GEX} indicates dealers are net long gamma: they buy when prices fall and sell when prices rise, \textit{dampening} volatility. \textbf{Negative GEX} indicates dealers are net short gamma: they sell when prices fall and buy when prices rise, \textit{amplifying} volatility through pro-cyclical hedging flows. The magnitude determines the scale of forced hedging.

This open interest-based approach (GEX\_OI) measures dealer inventory positioning rather than intraday flow. While volume-based GEX may better capture 0DTE dynamics, 30-day window aggregation smooths daily measurement noise, and the underlying hedging mechanism remains constant regardless of measurement method \citep{dim2025zero}.

\subsection{Multi-Phase Validation Strategy}

For regime detection, we employ a five-phase protocol spanning 2020--2025:

\textbf{Phase 1} (Q1 2024 Baseline, 52 windows): Establish baseline detection rate. Success criteria: detection substantially above chance while remaining selective---not universal matching, which would indicate base-rate guessing rather than structural discrimination.

\textbf{Phase 2} (Negative Controls, 809 windows): Three synthetic control types, each violating a specific regime criterion:
\begin{itemize}
    \item \textit{Shuffled} (404 windows): Randomize day order within real windows, destroying temporal structure while preserving aggregate statistics.
    \item \textit{Transitional} (202 windows): Generate windows with frequent sign flips ($>$8 per window), violating the stability criterion.
    \item \textit{Low-magnitude} (203 windows): Generate windows with magnitude below \$3B, violating the magnitude criterion.
\end{itemize}
Expected: $<$10\% false positive rate across all control types.

\textbf{Phase 3} (Full 2024, 223 windows): Test Q1 generalization across all 252 trading days.

\textbf{Phase 4} (2020 Comparison, 223 windows): Test pre-0DTE era (2020, $<$5\% 0DTE volume) against post-0DTE (2024, $\approx$48\%).

\textbf{Phase 5} (Multi-Year 2020--2025, 1,412 windows): Identify when structural market shift occurred.

For each LLM response, we extract stated metrics (persistence, magnitude, flips) and compare against ground truth. Windows with $>$5\% metric discrepancy are flagged for manual review.

\subsection{LLM Configuration}

We use OpenAI o4-mini \citep{openai2024reasoning} with temperature=1.0, max tokens=16,384, processed via Batch API (asynchronous, 100\% completion rate). The model receives a system message (``financial market analyst identifying persistent dealer gamma regimes''), a 30-day obfuscated GEX sequence with classification criteria, and outputs structured JSON with regime type, confidence (0--100), reasoning trace, and computed metrics. Total processing cost across all 2,221 evaluations was \$11.07.

\subsection{LLM Usage Disclosure}

In accordance with journal policy, we disclose the following use of AI tools: (1) OpenAI's o4-mini model serves as the \textit{subject of investigation}---the system whose structural reasoning capabilities are being validated; (2) Anthropic's Claude assisted with code development, data pipeline construction, and manuscript preparation. All model outputs were independently verified by the authors. All scientific interpretations and conclusions are solely those of the authors.
